\documentclass{article}
\usepackage[utf8]{vietnam}
\usepackage{times}
\usepackage{amsmath,amssymb}
\usepackage{tabularx}
\usepackage{array}
\usepackage{booktabs}
\usepackage{xcolor}
\usepackage[
  paperwidth=29.7cm,
  paperheight=16.706cm,
  top=0.9cm,
  bottom=0.9cm,
  left=1.2cm,
  right=1.2cm,
  footskip=0pt,
  headheight=0pt,
  headsep=0pt
]{geometry}

\pagestyle{empty}
\setlength{\parindent}{0pt}
\setlength{\parskip}{0pt}
\raggedbottom

\definecolor{brandRed}{RGB}{211, 47, 47}
\definecolor{brandGreen}{RGB}{46, 125, 50}
\definecolor{brandBlue}{RGB}{21, 101, 192}

\begin{document}\noindent\begin{minipage}[t][14.4cm][c]{\linewidth}
\begin{center}
{\fontsize{34}{44}\selectfont \textbf{\color{brandRed}LỚP KHTN THẦY THẠNH CÔ CHÂU}}\\[16pt]
{\fontsize{24}{32}\selectfont \textbf{CHUYÊN ĐỀ SINH HỌC 12}}\\[22pt]
{\fontsize{38}{48}\selectfont \textbf{\color{brandGreen}DỊCH MÃ VÀ MỐI QUAN HỆ TỪ GENE ĐẾN TÍNH TRẠNG}}\\[26pt]
{\fontsize{22}{30}\selectfont \textit{Hệ thống câu hỏi Trắc nghiệm nhiều lựa chọn, Đúng/Sai, Trả lời ngắn \& Tự luận}}\\[14pt]
{\fontsize{20}{28}\selectfont \textit{Tài liệu trình chiếu màn hình tương tác / Tivi 65 inch}}
\end{center}
\end{minipage}
\newpage


\begin{minipage}[t][14.4cm][t]{\linewidth}
\fontsize{22}{29}\selectfont

\textbf{Câu 1.} Phân tử nào sau đây được dùng làm khuôn cho quá trình dịch mã?\\[4pt]
\begin{tabularx}{\linewidth}{@{}X X@{}}
\textbf{A. } rRNA. & \textbf{B. } mRNA.\\[4pt]
\textbf{C. } tRNA. & \textbf{D. } Gene.
\end{tabularx}
\vspace*{\fill}
\textbf{Câu 2.} Khi nói về quá trình dịch mã, kết luận nào sau đây không đúng?\\[4pt]
\begin{tabularx}{\linewidth}{@{}X X@{}}
\textbf{A. } Bộ ba kết thúc quy định tổng hợp amino acid cuối cùng trên chuỗi polypeptide. & \textbf{B. } Trình tự các bộ ba trên mRNA quy định trình tự amino acid trong chuỗi polypeptide.\\[4pt]
\textbf{C. } Liên kết hydrogen được hình thành trước liên kết peptide. & \textbf{D. } Chiều chuyển dịch của ribosome trên mRNA là 5' $\rightarrow$ 3'.
\end{tabularx}

\end{minipage}
\newpage


\begin{minipage}[t][14.4cm][t]{\linewidth}
\fontsize{25}{33}\selectfont

\textbf{Câu 3.} Đối mã đặc hiệu trên phân tử tRNA được gọi là\\[4pt]
\begin{tabularx}{\linewidth}{@{}X X@{}}
\textbf{A. } codon. & \textbf{B. } amino acid.\\[4pt]
\textbf{C. } anticodon. & \textbf{D. } triplet.
\end{tabularx}
\vspace*{\fill}
\textbf{Câu 4.} Dịch mã thông tin di truyền trên bản mã sao thành trình tự amino acid trong chuỗi polypeptide là chức năng của\\[4pt]
\begin{tabularx}{\linewidth}{@{}X X@{}}
\textbf{A. } rRNA. & \textbf{B. } mRNA.\\[4pt]
\textbf{C. } tRNA. & \textbf{D. } RNA.
\end{tabularx}

\end{minipage}
\newpage


\begin{minipage}[t][14.4cm][t]{\linewidth}
\fontsize{20}{26}\selectfont

\textbf{Câu 5.} Các chuỗi polypeptide được tổng hợp trong tế bào nhân thực đều\\[4pt]
\begin{tabularx}{\linewidth}{@{}X X@{}}
\textbf{A. } kết thúc bằng Met. & \textbf{B. } bắt đầu bằng amino acid Met.\\[4pt]
\textbf{C. } bắt đầu bằng formyl-Met. & \textbf{D. } bắt đầu từ một phức hợp aa-tRNA.
\end{tabularx}
\vspace*{\fill}
\textbf{Câu 6.} Dịch mã là\\[4pt]
\begin{tabularx}{\linewidth}{@{}X X@{}}
\textbf{A. } quá trình tổng hợp RNA dựa trên trình tự nucleotide của DNA. & \textbf{B. } quá trình tổng hợp chuỗi polypeptide dựa trên trình tự nucleotide trên bản phiên mã của gene (mRNA).\\[4pt]
\textbf{C. } quá trình tạo 2 phân tử DNA mới giống như DNA ban đầu. & \textbf{D. } quá trình tổng hợp mRNA dựa trên trình tự nucleotide của chuỗi polypeptide.
\end{tabularx}

\end{minipage}
\newpage


\begin{minipage}[t][14.4cm][t]{\linewidth}
\fontsize{18}{23}\selectfont

\textbf{Câu 7.} Mã di truyền là\\[4pt]
\begin{tabularx}{\linewidth}{@{}X X@{}}
\textbf{A. } Trình tự nucleotide trên gene (DNA) quy định thành phần và trình tự amino acid trên phân tử protein và mRNA. & \textbf{B. } Trình tự nucleotide trên RNA quy định thành phần và trình tự amino acid trên phân tử protein, qua phân tử protein và mRNA.\\[4pt]
\textbf{C. } Trình tự nucleotide trên chuỗi polypeptide quy định thành phần và trình tự amino acid trên phân tử protein và mRNA. & \textbf{D. } Trình tự nucleotide trên phân tử mRNA quy định thành phần và trình tự amino acid trên phân tử protein.
\end{tabularx}
\vspace*{\fill}
\textbf{Câu 8.} Đâu là sơ đồ thể hiện mối quan hệ giữa gene và tính trạng?\\[4pt]
\begin{tabularx}{\linewidth}{@{}X X@{}}
\textbf{A. } Gene $\rightarrow$ mRNA $\rightarrow$ protein $\rightarrow$ tính trạng. & \textbf{B. } Protein $\rightarrow$ mRNA $\rightarrow$ gene $\rightarrow$ tính trạng.\\[4pt]
\textbf{C. } Tính trạng $\rightarrow$ protein $\rightarrow$ mRNA $\rightarrow$ gene. & \textbf{D. } Gene $\rightarrow$ protein $\rightarrow$ mRNA $\rightarrow$ tính trạng.
\end{tabularx}

\end{minipage}
\newpage


\begin{minipage}[t][14.4cm][t]{\linewidth}
\fontsize{25}{33}\selectfont

\textbf{Câu 9.} Trong quá trình dịch mã, bộ ba mã sao 3'AUC5' của mRNA khớp bổ sung với bộ ba đối mã nào sau đây?\\[4pt]
\begin{tabularx}{\linewidth}{@{}X X@{}}
\textbf{A. } 5'UAG3'. & \textbf{B. } 3'AUG5'.\\[4pt]
\textbf{C. } 3'UAG5'. & \textbf{D. } 3'UAC5'.
\end{tabularx}
\vspace*{\fill}
\textbf{Câu 10.} Từ 4 loại nucleotide, có thể tạo ra bao nhiêu loại codon?\\[4pt]
\begin{tabularx}{\linewidth}{@{}X X@{}}
\textbf{A. } 4. & \textbf{B. } 8.\\[4pt]
\textbf{C. } 32. & \textbf{D. } 64.
\end{tabularx}

\end{minipage}
\newpage


\begin{minipage}[t][14.4cm][t]{\linewidth}
\fontsize{20}{26}\selectfont

\textbf{Câu 11.} Quá trình dịch mã dừng lại khi\\[4pt]
\begin{tabularx}{\linewidth}{@{}X X@{}}
\textbf{A. } Ribosome tiếp xúc với vùng kết thúc nằm ở đầu 5' của mạch mã sao. & \textbf{B. } Ribosome tiếp xúc với vùng kết thúc nằm ở đầu 3' của mạch mã gốc.\\[4pt]
\textbf{C. } Ribosome tiếp xúc với bộ ba kết thúc trên mRNA. & \textbf{D. } Ribosome tiếp xúc với vùng kết thúc nằm ở đầu của mạch mã gốc.
\end{tabularx}
\vspace*{\fill}
\textbf{Câu 12.} Ngoài chức năng vận chuyển amino acid, RNA vận chuyển còn có chức năng quan trọng là\\[4pt]
\begin{tabularx}{\linewidth}{@{}X X@{}}
\textbf{A. } nhân tố trung gian vận chuyển thông tin di truyền từ nhân ra tế bào chất. & \textbf{B. } truyền thông tin di truyền qua các thế hệ cơ thể và thế hệ tế bào.\\[4pt]
\textbf{C. } cấu tạo nên ribosome là nơi xảy ra quá trình sinh tổng hợp protein. & \textbf{D. } nhận ra bộ ba mã sao tương ứng trên RNA thông tin theo nguyên tắc bổ sung.
\end{tabularx}

\end{minipage}
\newpage


\begin{minipage}[t][14.4cm][t]{\linewidth}
\fontsize{20}{26}\selectfont

\textbf{Câu 13.} Một phân tử mRNA chỉ chứa 3 loại nucleotide là A, U, C. Nhóm các bộ ba nào dưới đây có thể có trên mạch bổ sung của gene đã phiên mã ra mRNA nói trên?\\[4pt]
\begin{tabularx}{\linewidth}{@{}X X@{}}
\textbf{A. } TAG, GAA, ATA, ATG. & \textbf{B. } AAA, CCA, TAA, TCC.\\[4pt]
\textbf{C. } AAG, GTT, TCC, CAA. & \textbf{D. } ATC, TAG, GAC, GAA.
\end{tabularx}
\vspace*{\fill}
\textbf{Câu 14.} Cho biết các codon mã hóa các amino acid tương ứng: GGG – Gly; CCC – Pro; GCU – Ala; CGA – Arg; UCG – Ser; AGC – Ser. Một mạch gốc của một gene ở vi khuẩn có trình tự các nucleotide là 5'AGCCGACCCGGG3'. Nếu đoạn mạch gốc này mang thông tin mã hóa cho đoạn polypeptide có 4 amino acid thì trình tự của 4 amino acid đó là\\[4pt]
\begin{tabularx}{\linewidth}{@{}X X@{}}
\textbf{A. } Ser-Arg-Pro-Gly. & \textbf{B. } Ser-Ala-Gly-Pro.\\[4pt]
\textbf{C. } Pro-Gly-Ser-Ala. & \textbf{D. } Gly-Pro-Ser-Arg.
\end{tabularx}

\end{minipage}
\newpage


\begin{minipage}[t][14.4cm][t]{\linewidth}
\fontsize{22}{29}\selectfont

\textbf{Câu 15.} Khẳng định nào dưới đây không đúng khi nói về mã di truyền?\\[4pt]
\begin{tabularx}{\linewidth}{@{}X X@{}}
\textbf{A. } Mã di truyền là mã bộ ba. & \textbf{B. } Có nhiều mã di truyền cùng mang tín hiệu kết thúc dịch mã.\\[4pt]
\textbf{C. } Mỗi loại amino acid chỉ được xác định bởi một mã di truyền. & \textbf{D. } Một amino acid có thể được mã hóa bởi nhiều bộ ba.
\end{tabularx}
\vspace*{\fill}
\textbf{Câu 16.} Các codon nào dưới đây không mã hoá amino acid?\\[4pt]
\begin{tabularx}{\linewidth}{@{}X X@{}}
\textbf{A. } UAA, UAG, UGA. & \textbf{B. } AAU, GAU, UAC.\\[4pt]
\textbf{C. } AUA, UAA, UGC. & \textbf{D. } CUG, ACC, GUA.
\end{tabularx}

\end{minipage}
\newpage


\begin{minipage}[t][14.4cm][t]{\linewidth}
\fontsize{25}{33}\selectfont

\textbf{Câu 17.} Các mã di truyền khác nhau ở\\[4pt]
\begin{tabularx}{\linewidth}{@{}X X@{}}
\textbf{A. } số lượng và thành phần các nucleotide. & \textbf{B. } trình tự và thành phần nucleotide.\\[4pt]
\textbf{C. } thành phần các nucleotide. & \textbf{D. } số lượng và trình tự nucleotide.
\end{tabularx}
\vspace*{\fill}
\textbf{Câu 18.} Thành phần nào sau đây không tham gia quá trình dịch mã?\\[4pt]
\begin{tabularx}{\linewidth}{@{}X X@{}}
\textbf{A. } RNA. & \textbf{B. } Ribosome.\\[4pt]
\textbf{C. } DNA. & \textbf{D. } Amino acid.
\end{tabularx}

\end{minipage}
\newpage


\begin{minipage}[t][14.4cm][t]{\linewidth}
\fontsize{25}{33}\selectfont

\textbf{Câu 19.} Trong các bộ ba, có bao nhiêu bộ ba mã hóa cho các amino acid?\\[4pt]
\begin{tabularx}{\linewidth}{@{}X X@{}}
\textbf{A. } 64. & \textbf{B. } 63.\\[4pt]
\textbf{C. } 62. & \textbf{D. } 61.
\end{tabularx}
\vspace*{\fill}
\textbf{Câu 20.} Trong các bộ ba sau, bộ ba nào mã hóa cho amino acid Methionine?\\[4pt]
\begin{tabularx}{\linewidth}{@{}X X@{}}
\textbf{A. } UAA. & \textbf{B. } AUG.\\[4pt]
\textbf{C. } GGU. & \textbf{D. } CAA.
\end{tabularx}

\end{minipage}
\newpage


\begin{minipage}[t][14.4cm][t]{\linewidth}
\fontsize{18}{23}\selectfont

\textbf{Câu 21.} Cho các phát biểu sau:\\[3pt](1) Mã di truyền là thông tin về trình tự các amino acid được mã hóa dưới dạng trình tự các nucleotide trên mRNA.\\[3pt](2) Có 64 bộ ba được hình thành từ bốn loại nucleotide.\\[3pt](3) Số lượng, thành phần và trình tự sắp xếp các mã di truyền trên mRNA quy định số lượng, thành phần và trình tự sắp xếp các amino acid trên phân tử protein.\\[3pt](4) Ba bộ ba UAA, UGG và UCC đóng vai trò kết thúc dịch mã.\\[3pt]Trong các phát biểu trên, có bao nhiêu phát biểu đúng?\\[4pt]
\begin{tabularx}{\linewidth}{@{}X X@{}}
\textbf{A. } 1. & \textbf{B. } 2.\\[4pt]
\textbf{C. } 3. & \textbf{D. } 4.
\end{tabularx}
\vspace*{\fill}
\textbf{Câu 22.} Dịch mã là\\[4pt]
\begin{tabularx}{\linewidth}{@{}X X@{}}
\textbf{A. } quá trình tổng hợp chuỗi polypeptide thông qua thông tin di truyền được mã hóa trên phân tử mRNA. & \textbf{B. } quá trình tổng hợp chuỗi polynucleotide thông qua thông tin di truyền được mã hóa trên phân tử mRNA.\\[4pt]
\textbf{C. } quá trình tổng hợp chuỗi polypeptide thông qua thông tin di truyền được mã hóa trên phân tử DNA. & \textbf{D. } quá trình tổng hợp chuỗi polynucleotide thông qua thông tin di truyền được mã hóa trên phân tử DNA.
\end{tabularx}

\end{minipage}
\newpage


\begin{minipage}[t][14.4cm][t]{\linewidth}
\fontsize{20}{26}\selectfont

\textbf{Câu 23.} Trong các thành phần sau, có bao nhiêu thành phần trực tiếp tham gia vào quá trình tổng hợp chuỗi polypeptide?\\[3pt](1) Gene; (2) mRNA; (3) Amino acid; (4) tRNA; (5) Ribosome; (6) Enzyme; (7) rRNA; (8) RNA mồi.\\[4pt]
\begin{tabularx}{\linewidth}{@{}X X@{}}
\textbf{A. } 6. & \textbf{B. } 2.\\[4pt]
\textbf{C. } 3. & \textbf{D. } 4.
\end{tabularx}
\vspace*{\fill}
\textbf{Câu 24.} Cho các bộ ba mã hóa cho các amino acid tương ứng như sau: GGG – Gly; CCC – Pro; GCU – Ala; CGA – Arg; UCG – Ser; AGC – Ser. Đoạn mạch bổ sung của một gene ở vi khuẩn có trình tự các nucleotide: 3'-TCGGCTGGGGCCC-5'. Nếu đoạn mạch khuôn của gene này mang thông tin mã hóa cho đoạn polypeptide có 4 amino acid thì trình tự của 4 amino acid đó là\\[4pt]
\begin{tabularx}{\linewidth}{@{}X X@{}}
\textbf{A. } Ser - Ala - Gly - Pro. & \textbf{B. } Pro - Gly - Ser - Ala.\\[4pt]
\textbf{C. } Ser - Arg - Pro - Gly. & \textbf{D. } Gly - Pro - Ser - Arg.
\end{tabularx}

\end{minipage}
\newpage


\begin{minipage}[t][14.4cm][t]{\linewidth}
\fontsize{20}{26}\selectfont

\textbf{Câu 25.} Quá trình dịch mã dừng lại khi\\[4pt]
\begin{tabularx}{\linewidth}{@{}X X@{}}
\textbf{A. } ribosome tiếp xúc với vùng kết thúc nằm ở đầu 5' của mạch mã sao. & \textbf{B. } ribosome tiếp xúc với vùng kết thúc nằm ở đầu 3' của mạch mã gốc.\\[4pt]
\textbf{C. } ribosome tiếp xúc với bộ ba kết thúc trên mRNA. & \textbf{D. } ribosome tiếp xúc với vùng kết thúc nằm ở đầu 5' của mạch mã gốc.
\end{tabularx}
\vspace*{\fill}
\textbf{Câu 26.} Chọn từ/cụm từ thích hợp điền vào chỗ trống trong câu sau: Trình tự các nucleotide trên phân tử mRNA được ......... thành trình tự các amino acid trên phân tử protein.\\[4pt]
\begin{tabularx}{\linewidth}{@{}X X@{}}
\textbf{A. } dịch mã. & \textbf{B. } tái bản.\\[4pt]
\textbf{C. } phiên mã. & \textbf{D. } biểu hiện.
\end{tabularx}

\end{minipage}
\newpage


\begin{minipage}[t][14.4cm][t]{\linewidth}
\fontsize{20}{26}\selectfont

\textbf{Câu 27.} Các gene khác nhau có thể quy định các tính trạng khác nhau là do\\[4pt]
\begin{tabularx}{\linewidth}{@{}X X@{}}
\textbf{A. } các gene có kích thước khác nhau. & \textbf{B. } các gene có thể bị đột biến.\\[4pt]
\textbf{C. } các gene có vị trí khác nhau trong các tế bào khác nhau. & \textbf{D. } các gene có trình tự các nucleotide khác nhau.
\end{tabularx}
\vspace*{\fill}
\textbf{Câu 28.} Khi nói về dịch mã ở sinh vật nhân thực, nhận định nào sau đây là không đúng?\\[4pt]
\begin{tabularx}{\linewidth}{@{}X X@{}}
\textbf{A. } Trong cùng một thời điểm có thể có nhiều ribosome tham gia dịch mã trên một phân tử mRNA. & \textbf{B. } Amino acid mở đầu trong quá trình dịch mã là metionine.\\[4pt]
\textbf{C. } Khi dịch mã, ribosome dịch chuyển theo chiều 5' $\rightarrow$ 3' trên phân tử mRNA. & \textbf{D. } Khi dịch mã, ribosome dịch chuyển theo chiều 3' $\rightarrow$ 5' trên phân tử mRNA.
\end{tabularx}

\end{minipage}
\newpage


\begin{minipage}[t][14.4cm][t]{\linewidth}
\fontsize{18}{23}\selectfont

\textbf{Câu 29.} Khi nói về cơ chế dịch mã, có bao nhiêu nhận định không đúng trong các nhận định sau?\\[3pt](1) Trên một phân tử mRNA hoạt động của polysome giúp tạo ra các chuỗi polypeptide loại khác nhau.\\[3pt](2) Khi dịch mã, ribosome dịch chuyển theo chiều 3' – 5' trên phân tử mRNA.\\[3pt](3) Với bộ ba UAG trên mRNA thì tRNA mang đối mã là AUC.\\[3pt](4) Các chuỗi polypeptide sau dịch mã được cắt bỏ amino acid mở đầu và tiếp tục hình thành các bậc cấu trúc cao hơn để trở thành protein có hoạt tính sinh học.\\[3pt](5) Sau khi hoàn tất quá trình dịch mã, ribosome tách khỏi mRNA và giữ nguyên cấu trúc để chuẩn bị cho quá trình dịch mã tiếp theo.\\[3pt]Số nhận định không đúng là\\[4pt]
\begin{tabularx}{\linewidth}{@{}X X@{}}
\textbf{A. } 2. & \textbf{B. } 4.\\[4pt]
\textbf{C. } 1. & \textbf{D. } 3.
\end{tabularx}
\vspace*{\fill}
\textbf{Câu 30.} Phân tử nào sau đây mang bộ ba đối mã?\\[4pt]
\begin{tabularx}{\linewidth}{@{}X X@{}}
\textbf{A. } mRNA. & \textbf{B. } DNA.\\[4pt]
\textbf{C. } tRNA. & \textbf{D. } rRNA.
\end{tabularx}

\end{minipage}
\newpage


\begin{minipage}[t][14.4cm][t]{\linewidth}
\fontsize{18}{23}\selectfont

\textbf{Câu 31.} Ở tế bào nhân thực, quá trình dịch mã được bắt đầu bằng bộ ba nào trên mRNA và amino acid nào mở đầu chuỗi polypeptide?\\[4pt]
\begin{tabularx}{\linewidth}{@{}X X@{}}
\textbf{A. } AGU và formyl-Met. & \textbf{B. } AUG và formyl-Met.\\[4pt]
\textbf{C. } AUG và amino acid Met. & \textbf{D. } AGU và amino acid Met.
\end{tabularx}
\vspace*{\fill}
\textbf{Câu 32.} Trong giờ thực hành, giáo viên trình chiếu sơ đồ: DNA $\rightarrow$ (phiên mã) $\rightarrow$ mRNA $\rightarrow$ (dịch mã) $\rightarrow$ chuỗi polypeptide $\rightarrow$ protein $\rightarrow$ tính trạng. Ý nào sau đây mô tả chính xác vai trò của gene trong sơ đồ trên?\\[4pt]
\begin{tabularx}{\linewidth}{@{}X X@{}}
\textbf{A. } Gene trực tiếp tạo nên tính trạng mà không cần trung gian. & \textbf{B. } Gene quy định trình tự nucleotide trên mRNA; mRNA quy định trình tự amino acid trong protein; protein góp phần tạo nên tính trạng.\\[4pt]
\textbf{C. } Gene chỉ là vật chất lưu trữ, không tham gia quyết định tính trạng. & \textbf{D. } Gene chỉ tác động lên tính trạng khi môi trường thay đổi.
\end{tabularx}

\end{minipage}
\newpage


\begin{minipage}[t][14.4cm][t]{\linewidth}
\fontsize{25}{33}\selectfont

\textbf{Câu 33.} Quá trình dịch mã kết thúc khi\\[4pt]
\begin{tabularx}{\linewidth}{@{}X X@{}}
\textbf{A. } ribosome tiếp xúc với bộ ba AUG trên mRNA. & \textbf{B. } ribosome rời khỏi mRNA và trở về trạng thái tự do.\\[4pt]
\textbf{C. } ribosome tiếp xúc với một trong các bộ ba: UAA, UAG, UGA. & \textbf{D. } ribosome gắn methionine vào vị trí cuối cùng của chuỗi polypeptide.
\end{tabularx}
\vspace*{\fill}

\end{minipage}
\newpage


\begin{minipage}[t][14.4cm][t]{\linewidth}
\fontsize{22}{29}\selectfont
{\color{brandGreen}\textbf{PHẦN II: CÂU TRẮC NGHIỆM ĐÚNG / SAI}\\[8pt]}
\textbf{Câu 1.} Khi nói về đặc điểm của mã di truyền:
\vspace{8pt}

\begin{tabularx}{\linewidth}{@{}X r@{}}

\textbf{a) } Mã di truyền là thông tin về trình tự các axit amin được mã hoá dưới dạng trình tự các nuclêôtit trên mARN. & $[\quad]$\\[8pt]
\textbf{b) } Có 61 bộ ba mã hoá cho các axit amin và 3 bộ ba UAA, UAG, UGA đóng vai trò là mã kết thúc dịch mã. & $[\quad]$\\[8pt]
\textbf{c) } Bộ ba AUG chỉ mã hoá axit amin methionin và không đóng vai trò là mã mở đầu. & $[\quad]$\\[8pt]
\textbf{d) } Số lượng, thành phần và trật tự sắp xếp các bộ ba trên mARN quy định số lượng, thành phần và trật tự sắp xếp các amino acid trong chuỗi polypeptide. & $[\quad]$\\[8pt]
\end{tabularx}
\vfill
\end{minipage}
\newpage


\begin{minipage}[t][14.4cm][t]{\linewidth}
\fontsize{25}{33}\selectfont
{\color{brandGreen}\textbf{PHẦN II: CÂU TRẮC NGHIỆM ĐÚNG / SAI}\\[8pt]}
\textbf{Câu 2.} Khi xét về cơ sở phân tử và di truyền học:
\vspace{8pt}

\begin{tabularx}{\linewidth}{@{}X r@{}}

\textbf{a) } Axit nucleic là hợp chất hữu cơ có vai trò lưu giữ và truyền đạt thông tin di truyền. & $[\quad]$\\[8pt]
\textbf{b) } Nhiều mã di truyền cùng mã hoá một loại amino acid. & $[\quad]$\\[8pt]
\textbf{c) } Kiểu gen AaBb là kết quả của phép lai sau: P: AaBB × AaBB. & $[\quad]$\\[8pt]
\textbf{d) } Dòng thuần là tập hợp các cơ thể có kiểu hình giống nhau qua nhiều thế hệ tự thụ phấn hoặc giao phối cận huyết. & $[\quad]$\\[8pt]
\end{tabularx}
\vfill
\end{minipage}
\newpage


\begin{minipage}[t][14.4cm][t]{\linewidth}
\fontsize{25}{33}\selectfont
{\color{brandGreen}\textbf{PHẦN II: CÂU TRẮC NGHIỆM ĐÚNG / SAI}\\[8pt]}
\textbf{Câu 3.} Một đoạn phân tử DNA ở sinh vật nhân thực có trình tự nucleotide trên mạch gốc là: 3'...AAACAATGGGGA...5'.
\vspace{8pt}

\begin{tabularx}{\linewidth}{@{}X r@{}}

\textbf{a) } Trình tự nucleotide trên mạch bổ sung của đoạn DNA này là 5'...AAAGTTACCGGT...3'. & $[\quad]$\\[8pt]
\textbf{b) } Trình tự nucleotide trên mạch mRNA do gene trên phiên mã là 5'...UUUGUUACCCCU...3'. & $[\quad]$\\[8pt]
\textbf{c) } Chiều tổng hợp mạch mRNA theo chiều 3'......5'. & $[\quad]$\\[8pt]
\textbf{d) } Có 5 codon được tạo ra từ đoạn gene trên. & $[\quad]$\\[8pt]
\end{tabularx}
\vfill
\end{minipage}
\newpage


\begin{minipage}[t][14.4cm][t]{\linewidth}
\fontsize{19}{25}\selectfont
{\color{brandGreen}\textbf{PHẦN II: CÂU TRẮC NGHIỆM ĐÚNG / SAI}\\[8pt]}
\textbf{Câu 4.} Trong các vụ án hình sự hoặc xác định quan hệ huyết thống, các nhà khoa học thường sử dụng kỹ thuật "vân tay DNA". Chỉ cần một mẫu nhỏ như sợi tóc, giọt máu hoặc tế bào biểu mô, người ta có thể giải mã trình tự các nucleotide để xác định danh tính cá nhân với độ chính xác gần như tuyệt đối. Điều này dựa trên tính đặc trưng của hệ gene và các cơ sở phân tử của sự di truyền.\\[4pt]Dựa vào thông tin trên và kiến thức đã học, hãy xác định các nhận định sau là đúng hay sai?
\vspace{8pt}

\begin{tabularx}{\linewidth}{@{}X r@{}}

\textbf{a) } Kỹ thuật xác định huyết thống dựa trên nguyên tắc mỗi cá nhân có một hệ gene đặc trưng, trừ trường hợp sinh đôi cùng trứng. & $[\quad]$\\[8pt]
\textbf{b) } Mã di truyền là trình tự sắp xếp các amino acid trong phân tử protein, quy định các đặc điểm tính trạng của cơ thể. & $[\quad]$\\[8pt]
\textbf{c) } Chỉ từ 4 loại nucleotide (A, T, G, C), sự thay đổi trình tự sắp xếp và số lượng của chúng trên DNA đã tạo ra sự đa dạng vô tận của mã di truyền giữa các loài. & $[\quad]$\\[8pt]
\textbf{d) } Quá trình tái bản DNA diễn ra theo nguyên tắc bổ sung, giúp thông tin di truyền được truyền đạt nguyên vẹn từ tế bào mẹ sang các tế bào con. & $[\quad]$\\[8pt]
\end{tabularx}
\vfill
\end{minipage}
\newpage


\begin{minipage}[t][14.4cm][t]{\linewidth}
\fontsize{28}{38}\selectfont
{\color{brandGreen}\textbf{PHẦN III: CÂU TRẮC NGHIỆM TRẢ LỜI NGẮN}\\[10pt]}
\textbf{Câu 1.} Trên mạch mang mã gốc của gene có một bộ ba 3'AGC5'. Bộ ba tương ứng trên phân tử mRNA được phiên mã từ gene này là gì?
\vspace{20pt}

\textbf{Trả lời:} $\dots\dots\dots\dots\dots\dots\dots\dots\dots\dots\dots\dots\dots\dots\dots\dots\dots$
\vfill
\end{minipage}
\newpage


\begin{minipage}[t][14.4cm][t]{\linewidth}
\fontsize{28}{38}\selectfont
{\color{brandGreen}\textbf{PHẦN III: CÂU TRẮC NGHIỆM TRẢ LỜI NGẮN}\\[10pt]}
\textbf{Câu 2.} Phân tử nào trực tiếp làm khuôn cho quá trình dịch mã?
\vspace{20pt}

\textbf{Trả lời:} $\dots\dots\dots\dots\dots\dots\dots\dots\dots\dots\dots\dots\dots\dots\dots\dots\dots$
\vfill
\end{minipage}
\newpage


\begin{minipage}[t][14.4cm][t]{\linewidth}
\fontsize{28}{38}\selectfont
{\color{brandGreen}\textbf{PHẦN III: CÂU TRẮC NGHIỆM TRẢ LỜI NGẮN}\\[10pt]}
\textbf{Câu 3.} Ở sinh vật nhân thực, codon 5'AUG3' mã hóa loại amino acid nào?
\vspace{20pt}

\textbf{Trả lời:} $\dots\dots\dots\dots\dots\dots\dots\dots\dots\dots\dots\dots\dots\dots\dots\dots\dots$
\vfill
\end{minipage}
\newpage


\begin{minipage}[t][14.4cm][t]{\linewidth}
\fontsize{28}{38}\selectfont
{\color{brandGreen}\textbf{PHẦN III: CÂU TRẮC NGHIỆM TRẢ LỜI NGẮN}\\[10pt]}
\textbf{Câu 4.} Mạch khuôn của gene có đoạn 3' TATGGGCATGTA 5' thì mRNA được phiên mã từ mạch khuôn này có trình tự nucleotide như thế nào?
\vspace{20pt}

\textbf{Trả lời:} $\dots\dots\dots\dots\dots\dots\dots\dots\dots\dots\dots\dots\dots\dots\dots\dots\dots$
\vfill
\end{minipage}
\newpage


\begin{minipage}[t][14.4cm][t]{\linewidth}
\fontsize{28}{38}\selectfont
{\color{brandGreen}\textbf{PHẦN III: CÂU TRẮC NGHIỆM TRẢ LỜI NGẮN}\\[10pt]}
\textbf{Câu 5.} Giả sử có 3 loại nucleotide A, U, G thì phân tử mRNA có tối đa bao nhiêu loại mã di truyền mã hóa amino acid?
\vspace{20pt}

\textbf{Trả lời:} $\dots\dots\dots\dots\dots\dots\dots\dots\dots\dots\dots\dots\dots\dots\dots\dots\dots$
\vfill
\end{minipage}
\newpage


\begin{minipage}[t][14.4cm][t]{\linewidth}
\fontsize{25}{34}\selectfont
{\color{brandGreen}\textbf{PHẦN III: CÂU TRẮC NGHIỆM TRẢ LỜI NGẮN}\\[10pt]}
\textbf{Câu 6.} Một phân tử mRNA có 1200 nucleotide nhưng tính từ bộ ba mở đầu đến bộ ba kết thúc thì có tổng số 303 nucleotide. Phân tử mRNA này dịch mã có 6 ribosome trượt qua 1 lần đã cần môi trường cung cấp bao nhiêu amino acid?
\vspace{20pt}

\textbf{Trả lời:} $\dots\dots\dots\dots\dots\dots\dots\dots\dots\dots\dots\dots\dots\dots\dots\dots\dots$
\vfill
\end{minipage}
\newpage


\begin{minipage}[t][14.4cm][t]{\linewidth}
\fontsize{25}{34}\selectfont
{\color{brandGreen}\textbf{PHẦN III: CÂU TRẮC NGHIỆM TRẢ LỜI NGẮN}\\[10pt]}
\textbf{Câu 7.} Một phân tử mRNA có 1200 nucleotide nhưng tính từ bộ ba mở đầu đến bộ ba kết thúc thì có tổng số 606 nucleotide. Phân tử mRNA này dịch mã tạo ra chuỗi polypeptide. Tính số amino acid của mỗi chuỗi polypeptide hoàn chỉnh.
\vspace{20pt}

\textbf{Trả lời:} $\dots\dots\dots\dots\dots\dots\dots\dots\dots\dots\dots\dots\dots\dots\dots\dots\dots$
\vfill
\end{minipage}
\newpage


\begin{minipage}[t][14.4cm][t]{\linewidth}
\fontsize{28}{38}\selectfont
{\color{brandGreen}\textbf{PHẦN III: CÂU TRẮC NGHIỆM TRẢ LỜI NGẮN}\\[10pt]}
\textbf{Câu 8.} Quá trình dịch mã diễn ra ở đâu?
\vspace{20pt}

\textbf{Trả lời:} $\dots\dots\dots\dots\dots\dots\dots\dots\dots\dots\dots\dots\dots\dots\dots\dots\dots$
\vfill
\end{minipage}
\newpage


\begin{minipage}[t][14.4cm][t]{\linewidth}
\fontsize{28}{38}\selectfont
{\color{brandGreen}\textbf{PHẦN III: CÂU TRẮC NGHIỆM TRẢ LỜI NGẮN}\\[10pt]}
\textbf{Câu 9.} tRNA là gì?
\vspace{20pt}

\textbf{Trả lời:} $\dots\dots\dots\dots\dots\dots\dots\dots\dots\dots\dots\dots\dots\dots\dots\dots\dots$
\vfill
\end{minipage}
\newpage


\begin{minipage}[t][14.4cm][t]{\linewidth}
\fontsize{25}{34}\selectfont
{\color{brandGreen}\textbf{PHẦN III: CÂU TRẮC NGHIỆM TRẢ LỜI NGẮN}\\[10pt]}
\textbf{Câu 10.} Phân tử mRNA ở tế bào nhân sơ được phiên mã từ một gene có 3000 nucleotide sau đó tham gia vào quá trình dịch mã. Quá trình tổng hợp protein có 5 ribosome cùng trượt trên mRNA đó. Số amino acid môi trường cần cung cấp để hoàn tất quá trình dịch mã trên là bao nhiêu?
\vspace{20pt}

\textbf{Trả lời:} $\dots\dots\dots\dots\dots\dots\dots\dots\dots\dots\dots\dots\dots\dots\dots\dots\dots$
\vfill
\end{minipage}
\newpage


\begin{minipage}[t][14.4cm][t]{\linewidth}
\fontsize{22}{30}\selectfont
{\color{brandGreen}\textbf{PHẦN III: CÂU TRẮC NGHIỆM TRẢ LỜI NGẮN}\\[10pt]}
\textbf{Câu 11.} Các bộ ba trên mRNA mã hóa các amino acid tương ứng sau: 5'UUG3' – Leu, 5'AUC3' – Ile, 5'ACC3' – Thr, 5'AAU3' – Asn, 5'AUG3' – Met (amino acid mở đầu), 5'GUA3' – Val, 5'UGC3' – Cys, 5'UAG 3' – Kết thúc, 5'GGC3' – Gly. Một phân tử mRNA ngắn của vi khuẩn có trình tự các nucleotide như sau: 5'AUG AUC UUG ACC AUC AAU UGC GGC UGC UGA3'. Xác định trình tự nucleotide trên mạch bổ sung của gene.
\vspace{20pt}

\textbf{Trả lời:} $\dots\dots\dots\dots\dots\dots\dots\dots\dots\dots\dots\dots\dots\dots\dots\dots\dots$
\vfill
\end{minipage}
\newpage


\begin{minipage}[t][14.4cm][t]{\linewidth}
\fontsize{22}{30}\selectfont
{\color{brandGreen}\textbf{PHẦN III: CÂU TRẮC NGHIỆM TRẢ LỜI NGẮN}\\[10pt]}
\textbf{Câu 12.} Cho biết các codon mã hóa các amino acid tương ứng như sau: 5'GGG3' – Gly; 5'CCC3' – Pro; 5'GCU3' – Ala; 5'CGA3' – Arg; 5'UCG3' – Ser; 5'AGC3' – Ser; 5'UAC3' – Tyr. Một đoạn mạch gốc của một gene ở vi khuẩn mang thông tin mã hóa cho đoạn polypeptide có 5 amino acid có trình tự các nucleotide là 3'CCC-AGC-ATG-CGA-GGG5'. Xác định trình tự của 5 amino acid do đoạn gene này quy định tổng hợp.
\vspace{20pt}

\textbf{Trả lời:} $\dots\dots\dots\dots\dots\dots\dots\dots\dots\dots\dots\dots\dots\dots\dots\dots\dots$
\vfill
\end{minipage}
\newpage


\begin{minipage}[t][14.4cm][t]{\linewidth}
\fontsize{28}{38}\selectfont
{\color{brandGreen}\textbf{PHẦN III: CÂU TRẮC NGHIỆM TRẢ LỜI NGẮN}\\[10pt]}
\textbf{Câu 13.} Một bộ ba ribonucleotide liền kề trên mRNA quy định một amino acid (hoặc mang tín hiệu kết thúc) được gọi là gì?
\vspace{20pt}

\textbf{Trả lời:} $\dots\dots\dots\dots\dots\dots\dots\dots\dots\dots\dots\dots\dots\dots\dots\dots\dots$
\vfill
\end{minipage}
\newpage


\begin{minipage}[t][14.4cm][t]{\linewidth}
\fontsize{28}{38}\selectfont
{\color{brandGreen}\textbf{PHẦN III: CÂU TRẮC NGHIỆM TRẢ LỜI NGẮN}\\[10pt]}
\textbf{Câu 14.} Từ 4 loại ribonucleotide (A, U, G, C), có thể hình thành tất cả bao nhiêu loại codon?
\vspace{20pt}

\textbf{Trả lời:} $\dots\dots\dots\dots\dots\dots\dots\dots\dots\dots\dots\dots\dots\dots\dots\dots\dots$
\vfill
\end{minipage}
\newpage


\begin{minipage}[t][14.4cm][t]{\linewidth}
\fontsize{28}{38}\selectfont
{\color{brandGreen}\textbf{PHẦN III: CÂU TRẮC NGHIỆM TRẢ LỜI NGẮN}\\[10pt]}
\textbf{Câu 15.} Trong tổng số 64 bộ ba mã di truyền, có bao nhiêu bộ ba mã hóa cho các amino acid?
\vspace{20pt}

\textbf{Trả lời:} $\dots\dots\dots\dots\dots\dots\dots\dots\dots\dots\dots\dots\dots\dots\dots\dots\dots$
\vfill
\end{minipage}
\newpage


\begin{minipage}[t][14.4cm][t]{\linewidth}
\fontsize{28}{38}\selectfont
{\color{brandGreen}\textbf{PHẦN III: CÂU TRẮC NGHIỆM TRẢ LỜI NGẮN}\\[10pt]}
\textbf{Câu 16.} Bộ ba làm nhiệm vụ khởi đầu quá trình dịch mã trên mRNA là bộ ba nào?
\vspace{20pt}

\textbf{Trả lời:} $\dots\dots\dots\dots\dots\dots\dots\dots\dots\dots\dots\dots\dots\dots\dots\dots\dots$
\vfill
\end{minipage}
\newpage


\begin{minipage}[t][14.4cm][t]{\linewidth}
\fontsize{28}{38}\selectfont
{\color{brandGreen}\textbf{PHẦN III: CÂU TRẮC NGHIỆM TRẢ LỜI NGẮN}\\[10pt]}
\textbf{Câu 17.} Khi gặp một trong ba bộ ba nào (viết theo chiều 5' $\rightarrow$ 3') trên phân tử mRNA thì quá trình dịch mã sẽ dừng lại?
\vspace{20pt}

\textbf{Trả lời:} $\dots\dots\dots\dots\dots\dots\dots\dots\dots\dots\dots\dots\dots\dots\dots\dots\dots$
\vfill
\end{minipage}
\newpage


\begin{minipage}[t][14.4cm][t]{\linewidth}
\fontsize{28}{38}\selectfont
{\color{brandGreen}\textbf{PHẦN III: CÂU TRẮC NGHIỆM TRẢ LỜI NGẮN}\\[10pt]}
\textbf{Câu 18.} Bào quan nào trong tế bào đóng vai trò là nơi trực tiếp diễn ra quá trình dịch mã tổng hợp chuỗi polypeptide?
\vspace{20pt}

\textbf{Trả lời:} $\dots\dots\dots\dots\dots\dots\dots\dots\dots\dots\dots\dots\dots\dots\dots\dots\dots$
\vfill
\end{minipage}
\newpage


\begin{minipage}[t][14.4cm][t]{\linewidth}
\fontsize{28}{38}\selectfont
{\color{brandGreen}\textbf{PHẦN III: CÂU TRẮC NGHIỆM TRẢ LỜI NGẮN}\\[10pt]}
\textbf{Câu 19.} Phân tử RNA nào giữ vai trò mang bộ ba đối mã (anticodon) và trực tiếp vận chuyển amino acid tương ứng đến ribosome?
\vspace{20pt}

\textbf{Trả lời:} $\dots\dots\dots\dots\dots\dots\dots\dots\dots\dots\dots\dots\dots\dots\dots\dots\dots$
\vfill
\end{minipage}
\newpage


\begin{minipage}[t][14.4cm][t]{\linewidth}
\fontsize{28}{38}\selectfont
{\color{brandGreen}\textbf{PHẦN III: CÂU TRẮC NGHIỆM TRẢ LỜI NGẮN}\\[10pt]}
\textbf{Câu 20.} Quá trình dịch mã tổng hợp chuỗi polypeptide được chia làm mấy giai đoạn chính?
\vspace{20pt}

\textbf{Trả lời:} $\dots\dots\dots\dots\dots\dots\dots\dots\dots\dots\dots\dots\dots\dots\dots\dots\dots$
\vfill
\end{minipage}
\newpage


\begin{minipage}[t][14.4cm][t]{\linewidth}
\fontsize{26}{35}\selectfont
{\color{brandGreen}\textbf{PHẦN IV: CÂU HỎI TỰ LUẬN}\\[10pt]}
\textbf{Câu 1.} Một phân tử mRNA có tổng cộng 600 nucleotide (kể cả bộ ba kết thúc). Chuỗi polypeptide hoàn chỉnh được tổng hợp có bao nhiêu amino acid?
\vspace{14pt}

\textbf{Bài làm:}
\vfill
\end{minipage}
\newpage


\begin{minipage}[t][14.4cm][t]{\linewidth}
\fontsize{26}{35}\selectfont
{\color{brandGreen}\textbf{PHẦN IV: CÂU HỎI TỰ LUẬN}\\[10pt]}
\textbf{Câu 2.} Để tổng hợp một chuỗi polypeptide gồm 399 amino acid, số lượng nucleotide tối thiểu cần có trên phân tử mRNA (bao gồm codon kết thúc) là bao nhiêu nucleotide?
\vspace{14pt}

\textbf{Bài làm:}
\vfill
\end{minipage}
\newpage


\begin{minipage}[t][14.4cm][t]{\linewidth}
\fontsize{23}{31}\selectfont
{\color{brandGreen}\textbf{PHẦN IV: CÂU HỎI TỰ LUẬN}\\[10pt]}
\textbf{Câu 3.} Một đoạn phân tử mRNA có trình tự như sau:\\[6pt]5' – AUGGCGAAGCCCUGCGAGUAA – 3'\\[6pt]a. Đoạn phân tử mRNA trên chứa bao nhiêu mã bộ ba?\\[6pt]b. Chuỗi polypeptide được dịch mã từ đoạn phân tử mRNA trên có bao nhiêu amino acid?
\vspace{14pt}

\textbf{Bài làm:}
\vfill
\end{minipage}
\newpage


\begin{minipage}[t][14.4cm][t]{\linewidth}
\fontsize{23}{31}\selectfont
{\color{brandGreen}\textbf{PHẦN IV: CÂU HỎI TỰ LUẬN}\\[10pt]}
\textbf{Câu 4.} Một nhà khoa học tổng hợp phân tử mRNA nhân tạo với vùng mã hoá protein chứa 1200 nucleotide (bao gồm cả mã mở đầu và mã kết thúc). Nhà khoa học thực hiện phản ứng dịch mã phân tử mRNA mới tổng hợp trong tế bào vi khuẩn E. coli. Hãy cho biết chuỗi polypeptide được dịch mã có bao nhiêu amino acid?
\vspace{14pt}

\textbf{Bài làm:}
\vfill
\end{minipage}
\newpage


\begin{minipage}[t][14.4cm][t]{\linewidth}
\fontsize{23}{31}\selectfont
{\color{brandGreen}\textbf{PHẦN IV: CÂU HỎI TỰ LUẬN}\\[10pt]}
\textbf{Câu 5.} Cho sơ đồ sau: Gene (DNA) ——(1)$\longrightarrow$ mRNA ——(2)$\longrightarrow$ Protein $\longrightarrow$ Tính trạng\\[6pt]a) Nêu tên và sản phẩm của quá trình 1, quá trình 2.\\[6pt]b) Giải thích mối quan hệ giữa gene (DNA), mRNA, protein và tính trạng.
\vspace{14pt}

\textbf{Bài làm:}
\vfill
\end{minipage}
\newpage


\begin{minipage}[t][14.4cm][t]{\linewidth}
\fontsize{26}{35}\selectfont
{\color{brandGreen}\textbf{PHẦN IV: CÂU HỎI TỰ LUẬN}\\[10pt]}
\textbf{Câu 6.} Cho đoạn phân tử mRNA có trình tự sắp xếp các nucleotide như sau:\\[6pt]5'...A–C–C–U–C–A–U–U–A–G–C–A–G–C–G...3'\\[6pt]Viết trình tự nucleotide của cả đoạn gene phiên mã đoạn mRNA trên?
\vspace{14pt}

\textbf{Bài làm:}
\vfill
\end{minipage}
\newpage


\begin{minipage}[t][14.4cm][t]{\linewidth}
\fontsize{23}{31}\selectfont
{\color{brandGreen}\textbf{PHẦN IV: CÂU HỎI TỰ LUẬN}\\[10pt]}
\textbf{Câu 7.} Một đoạn mạch của gene có cấu trúc như sau:\\[6pt]Mạch 1: – A – T – G – C – T – C – G –\\[6pt]Mạch 2: – T – A – C – G – A – G – C –\\[6pt]Xác định trình tự các đơn phân của đoạn mạch RNA được tổng hợp từ mạch 2.
\vspace{14pt}

\textbf{Bài làm:}
\vfill
\end{minipage}
\newpage


\begin{minipage}[t][14.4cm][t]{\linewidth}
\fontsize{23}{31}\selectfont
{\color{brandGreen}\textbf{PHẦN IV: CÂU HỎI TỰ LUẬN}\\[10pt]}
\textbf{Câu 8.} Cho biết các bộ ba trên mRNA mã hóa cho các amino acid tương ứng là: 5'CUG3' - Leu, 5'GUC3' - Val, 5'ACG3' - Thr, 5'GCA3' - Ala. Từ đoạn mạch gốc chứa 4 mã di truyền của một gene có trình tự đơn phân 5'CAGCGTGACCAG3' phiên mã tổng hợp đoạn mRNA. Theo cơ chế dịch mã, hãy viết đoạn polypeptide được tạo ra từ đoạn mRNA trên.
\vspace{14pt}

\textbf{Bài làm:}
\vfill
\end{minipage}
\newpage


\begin{minipage}[t][14.4cm][t]{\linewidth}
\fontsize{19}{25}\selectfont
\begin{center}
{\fontsize{26}{34}\selectfont \textbf{\color{brandRed}BẢNG ĐÁP ÁN THAM KHẢO}}\\[10pt]
\textbf{\color{brandBlue}PHẦN I: TRẮC NGHIỆM NHIỀU PHƯƠNG ÁN LỰA CHỌN}\\[4pt]
\begin{tabular}{|c|c||c|c||c|c||c|c||c|c||c|c|}
\hline
\textbf{Câu} & \textbf{ĐA} & \textbf{Câu} & \textbf{ĐA} & \textbf{Câu} & \textbf{ĐA} & \textbf{Câu} & \textbf{ĐA} & \textbf{Câu} & \textbf{ĐA} & \textbf{Câu} & \textbf{ĐA}\\
\hline
1 & B & 7 & D & 13 & B & 19 & D & 25 & C & 31 & C\\
2 & A & 8 & A & 14 & C & 20 & B & 26 & A & 32 & B\\
3 & C & 9 & A & 15 & C & 21 & B & 27 & D & 33 & C\\
4 & C & 10 & D & 16 & A & 22 & A & 28 & D & & \\
5 & B & 11 & C & 17 & B & 23 & A & 29 & B & & \\
6 & B & 12 & D & 18 & C & 24 & D & 30 & C & & \\
\hline
\end{tabular}

\vspace{10pt}
\textbf{\color{brandBlue}PHẦN II: TRẮC NGHIỆM ĐÚNG / SAI}\\[4pt]
\begin{tabular}{|c|c|c|c|c|}
\hline
\textbf{Câu} & \textbf{Ý a} & \textbf{Ý b} & \textbf{Ý c} & \textbf{Ý d}\\
\hline
\textbf{1} & Đúng & Đúng & Sai & Đúng\\
\textbf{2} & Đúng & Đúng & Sai & Sai\\
\textbf{3} & Sai & Đúng & Sai & Sai\\
\textbf{4} & Đúng & Sai & Đúng & Đúng\\
\hline
\end{tabular}

\vspace{10pt}
\textbf{\color{brandBlue}PHẦN III: TRẢ LỜI NGẮN}\\[4pt]
\begin{tabular}{|c|c||c|c||c|c||c|c|}
\hline
\textbf{Câu} & \textbf{Đáp án} & \textbf{Câu} & \textbf{Đáp án} & \textbf{Câu} & \textbf{Đáp án} & \textbf{Câu} & \textbf{Đáp án}\\
\hline
1 & 5'UCG3' & 6 & 600 & 11 & 5'ATG...TGA3' & 16 & 5'AUG3'\\
2 & mRNA & 7 & 200 & 12 & Gly-Ser-Tyr-Ala-Pro & 17 & UAA, UAG, UGA\\
3 & Met & 8 & Tế bào chất & 13 & Codon & 18 & Ribosome\\
4 & 5'AUA...U3' & 9 & tRNA & 14 & 64 & 19 & tRNA\\
5 & 24 & 10 & 2495 & 15 & 61 & 20 & 3 giai đoạn\\
\hline
\end{tabular}
\end{center}
\end{minipage}
\newpage

\end{document}